% Clase 27 --- Interferencia de dos fuentes en el plano (§2.3--2.4).
% Los circulos son las crestas de cada fuente. Donde se cruza cresta con cresta
% hay refuerzo; entre medio quedan lineas donde siempre se cancelan, y esas
% lineas son hiperbolas (diferencia de caminos constante).
\begin{center}
\begin{tikzpicture}[scale=1]

  \coordinate (S1) at (-1.1,0);
  \coordinate (S2) at (1.1,0);

  % crestas de cada fuente, separadas lambda
  \foreach \rr in {1.1,2.2,3.3,4.4}{
    \draw[figaux] (S1) circle (\rr);
    \draw[figgray!45, dashed, line width=0.6pt] (S2) circle (\rr);}

  \filldraw[figamber] (S1) circle (2.4pt);
  \filldraw[figamber] (S2) circle (2.4pt);
  \node[figlbl, figamber, anchor=north, fill=white, inner sep=1.5pt]
    at (-1.1,-0.16) {$S_1$};
  \node[figlbl, figamber, anchor=north, fill=white, inner sep=1.5pt]
    at (1.1,-0.16) {$S_2$};

  % lineas nodales: hiperbolas |r1 - r2| = (n+1/2) lambda, con c = 1.1
  \foreach \aa/\bb in {0.275/1.0651, 0.825/0.7276}{
    \draw[figred, line width=0.9pt, smooth, samples=90,
          variable=\pp, domain=-1.75:1.75]
      plot ({\aa*cosh(\pp)}, {\bb*sinh(\pp)});
    \draw[figred, line width=0.9pt, smooth, samples=90,
          variable=\pp, domain=-1.75:1.75]
      plot ({-\aa*cosh(\pp)}, {\bb*sinh(\pp)});}

  % la linea central: constructiva siempre
  \draw[figblue, line width=1.1pt] (0,-3.4) -- (0,3.4);
  \node[figlbl, figblue, anchor=west, fill=white, inner sep=1.5pt]
    at (0.1,3.05) {refuerzo: $r_1 = r_2$};
  \node[figlbl, figred, anchor=west, fill=white, inner sep=1.5pt]
    at (2.0,2.2) {líneas oscuras: $|r_1 - r_2| = (n+\tfrac12)\lambda$};

\end{tikzpicture}\\[0.5em]
{\small\itshape Cada familia de círculos marca las crestas de una fuente,
separadas $\lambda$. Donde una cresta se cruza con otra cresta ---o un valle con
un valle--- las dos ondas se suman y hay \textbf{interferencia constructiva};
donde una cresta cae sobre un valle, se cancelan. Lo que ordena el dibujo es que
esas condiciones no dependen del tiempo sino sólo de la \textbf{diferencia de
caminos} $r_1 - r_2$: por eso los puntos de refuerzo no están sueltos sino
alineados sobre curvas fijas, alternadas con las curvas oscuras. La mediatriz es
siempre constructiva porque ahí las dos fuentes están a igual distancia. Escrito
en términos de fase, hay refuerzo cuando $\phi_2 - \phi_1 = 2\pi n$ y
cancelación cuando vale $\pi + 2\pi n$; lo que importa es la \textbf{diferencia}
de fases y no cada fase por separado ---conviene recordar que $-A\sen(kx)$ es lo
mismo que $A\sen(kx+\pi)$, de modo que un signo menos ya es medio ciclo de
desfasaje. Todo esto exige que las dos fuentes mantengan su desfasaje en el
tiempo, es decir, que sean coherentes.}
\end{center}
